# E8 Spectral Scan & Analysis Software PRD

**Goal:** Perform a **Blind Spectral Analysis** of the $E\_8$ lattice configuration space defined in Paper II. Generate the spectrum of allowed persistent states, filter them, and identify where they cluster relative to experimental data.

# Phase 1: The Physics Engine

**Reference:** Paper II, Sections II (*Configuration Space*) & III (*Persistence Filters*).

* Requirement: Implement the Generative Equation (Section II.C) and the 5 Hard Filters (Section III.A) exactly as defined in the text.  
* Constants: Use the derived values from Paper I ($\\alpha^{-1} \= 137.035999$, etc.).

**Add Function:** `generate_neutrino_spectrum()`

* Logic: Implement the inverse formula: $M\_n \= M\_{scale} / \\Delta^n$.  
* Generations: Iterate $n \\in {2, 3, 4}$ (corresponding to $\\nu\_3, \\nu\_2, \\nu\_1$).  
* Constants: $M\_{scale} \= m\_e \\cdot (\\pi\\alpha / (\\nu\\Delta))$.  
* Refractive Index: Apply $\\eta\_W \= 6$ to calculate the "Weak/Oscillation" mass ($M\_{weak} \= M\_{rest} \\cdot \\eta\_W$).

# Phase 2: The Landscape Analysis

1. Cluster Detection: Sort survivors. Cluster values within 0.15% tolerance (Isospin Limit).  
2. Forensic Mapping: Match cluster centers to `PARTICLE_DB`.  
   * Metric: Relative Residual ($R$) $= |M\_{calc} \- M\_{exp}| / M\_{exp}$.  
3. Stability Flagging:  
   * Persistence Margin (PM): $3.5 \\times 10^{-6} m\_e$.  
   * Status \= LOCKED if $R \\le PM$.  
   * Status \= UNSTABLE if $R \> PM$.  
4. Precision Tests: Implement `perform_precision_tests()`:  
   * Test 1: Kaon Splitting:  
     * Calculate $M(K^0)$ and $M(K^+)$ using the exact formulas from Paper II (Isospin Drag).  
     * Formula: $\\Delta M\_K \= D(\\sigma \- \\pi) \+ 1/D$.  
     * Compare to experimental $3.934$ MeV ($7.68 m\_e$).  
   * Test 2: Proton Radius:  
     * Calculate $r\_p \= D \\cdot \\hbar / (m\_p c)$.  
     * Compare to CODATA ($0.8414$ fm).  
   * Output: A "Pass/Fail" report block with % accuracy.

## CLI Arguments:

* `--export_csv`: Saves the full dataset to `E8_spectral_scan.csv`.  
* `--particle [NAME]`: filters output to show only formulas for a specific particle (e.g., "Proton") for detailed inspection.  
* `--summary`: Outputs the "Census Report" (Counts per Tier, Total Candidates, Locked Count).

* Add function: perform\_precision\_tests()  
  * Test 1: Kaon Splitting:  
    * Calculate $M(K^0)$ and $M(K^+)$ using the exact formulas from Paper II (Isospin Drag).  
    * Compute $\\Delta M$ and compare to experimental $3.934$ MeV.  
  * Test 2: Proton Radius:  
    * Calculate $r\_p \= 4 \\cdot \\hbar / (m\_p c)$.  
    * Compare to CODATA ($0.8414$ fm).  
  * Output: A "Pass/Fail" report block with % accuracy.

# Phase 3: Output Requirements

1\. The Standard Model Map (Table)

* Goal: Verify the 12 Fermions.  
* Columns: Name | Tier | Formula | Mass ($m\_e$) | Residual | Stability | Count  
* Checks:  
  * Did the Neutron appear via the Half-Step ($k=0.5$)?  
  * Did the Bottom appear via the Saturated Scalar ($S\_b$)?

2\. The Extended Hadron Map (Table)

* Goal: Map the Mesons/Baryons in the DB.  
* Cols: Same as above.

3\. The Boson Exclusion Report

* Goal: Verify Linear Search fails to find Non-Linear Bosons.  
* Check for W, Z, Higgs.  
* Pass Condition: They are MISSING or have high residuals ($\>1%$).  
* *Output Note:* "Non-detection of linear Boson resonances confirms non-linear structure ($Spin \\ge 1$)."

4\. The Neutrino Spectrum

* Call `generate_neutrino_spectrum()`.  
* Output: Print a table comparing Rest Mass (Cosmology) vs Weak Mass (Oscillation) for all 3 generations.

5\. Precision Validation

* Call `perform_precision_tests()` and output results.

6\. The Ghost Report

* Goal: Identify high-density solutions ($N\_{count} \> 50$) that do *not* match the DB.  
* Flag: Any ghost with $R \\le PM$ (Dark Matter candidates).

7\. The Structural Component Check (The "LEGO" Test) Search the raw spectrum for these exact mass values (Geometric Anchors), even if they aren't particles. Output FOUND/MISSING.

* The Volumetric Neutral Load:  
  * Target Mass: $\\Delta^3 / D \= 79,507 / 4 \\approx \\mathbf{19,877 \\ m\_e}$.  
  * *Expectation:* MISSING (Validates it is a load, not a free particle).  
* The Manifold Anchors (Quark Bases):  
  * Up Basis: Target $\\approx 4.0 m\_e$ ($D$).  
  * Down Basis: Target $\\approx 9.0 m\_e$ ($D+\\sigma$).  
  * *Expectation:* FOUND (Validates the scalar anchor inputs).

8\. The Decay Power Law Data & Plots

4\. Power Law Data Generation:

* Scope Restriction: Include only Fundamental Particles (`Lepton`, `Quark`, `Boson`). Exclude composite Hadrons (Mesons, Baryons) as their decay widths are dominated by binding energy dynamics which obscure the geometric residual signal in a linear scan.  
* Filter: Identify unstable particles where $R \> PM$.  
* Distillation: For each unique species, select the Single Best Candidate (lowest $|R|$).  
* Exclusion: Strictly exclude meta-stable particles ($\\tau \> 10^{-4} s$) to isolate geometric instability from topological protection.  
* Output: A clean list of unique tuples `(Name, Best_Residual, Lifetime)` for regression. This prevents statistical artifacts caused by multiple candidate solutions for a single mass.  
* Analysis: Calculate Slope ($\\beta$) and Correlation ($r^2$) programmatically using `scipy.stats.linregress`.  
* Unit Test: If $\\beta$ is between $-2.0$ and $-7.0$ and $R^2 \> 0.8$, print `[CONFIRMED] Geometric Phase Space (Slope Negative, High Correlation)`.  
* *Note:* The relaxed slope range accounts for the "Coarse Scan" limitation; the simplified generator cannot find the perfect deep-decimal residuals for light particles, which artificially steepens the slope, but the negative correlation must remain strong.  
* Output the full list

* Visualization:  
* Plot 1: The Discovery Landscape (`spectral_landscape.png`)  
* Title: “Geometric Resonance Spectrum: The Emergence of the Standard Model”  
* Method: Use a Histogram or Kernel Density Estimation (KDE) on the log-mass axis to show "Mountains of Convergence." across the four:  
* Color Schema:  
  * Background (Noise): Light Grey (Low opacity).  
  * Standard Model (Fermions): Electric Blue (High opacity, labeled).  
  * Hadrons/Mesons: Forest Green (Medium opacity).  
  * Ghosts: Crimson Red (For high-density unknown peaks).  
  * Fade the Noise: Make the Grey (Noise) and Red (Ghosts) much thinner and more transparent.  
  * Boost the Signal: Make the Blue (SM) and Green (Hadrons) lines thicker and the markers larger.  
  * Clean the Annotations: Move text slightly to avoid overlapping lines and ensure the font sizes are consistent.  
  * Gridlines: Add faint vertical gridlines to help the eye align mass values.  
* Annotations: Implemented a best\_peaks filter to identify the "Main" spectral line for every particle found and labeling e , μ , π , K , p , n , Δ , J / ψ , Υ , W , Z , t e,μ,π,K,p,n,Δ,J/ψ,Υ,W,Z,t clearly at their strongest resonance points.  
* X-Axis: Log scale ($m\_e$).  
* Y-Axis: Log scale (Solution Density).


* Plot 2: The Decay Power Law (`decay_regression.png`)  
* Data: Plot `log10(Residual)` vs `log10(Lifetime)`.  
* Analysis: Perform a Linear Regression (`scipy.stats.linregress`).  
* Visual: Scatter plot points \+ The Regression Line.  
* Text on Plot: Display the equation $y \= mx \+ b$, specifically the slope $\\beta$ (Target $\\approx \-3.0$) and $R^2$.  
* Constraint: Plot only the Fundamental Particles used in the regression (exclude Hadrons/Stable particles from the chart to reduce noise).  
* Annotation: Label each point with the particle name (e.g., "Muon", "Top").

# Coding Guidelines

* Maximize performance in data analysis, replace nested loops with running state variables and move static calculations outside of loops  
* When outputting tables, make a markdown table  
* Include the most persistent formula when outputting mapping  
* Prefer to use more functional style code  
* Math: Use the updated operator values:  
  * $\\alpha^{-1} \= 137.035999$  
  * $\\Delta \= 43, \\sigma \= 5, \\nu \= 16, \\chi \= 2, D \= 4$.  
* Structure: Separate files  
  * physics\_constants.py (DELTA, NU, SIGMA, D, CHI, H\_SYS, H\_FULL, ALPHA\_INV, MASS\_GAP, VEV\_GEV, ME\_GEV, VEV\_ME, S\_B, S\_T, PERSISTENCE\_MARGIN, and PARTICLE\_DB)  
  * generator.py generate\_spectrum(verbose=True)  
  * analyzer.py: analyze\_spectrum(survivors), perform\_structural\_checks()  
  * plot.py with 2 functions for the 2 visualizations plot\_spectral\_landscape(clusters, mapped, ghosts) and plot\_decay\_regression(clean\_data, slope, r2).  
  * main.py

// TODO Populate correctly with full PDG catalog and re-run

\# EXPANDED PARTICLE DATABASE (Units: Electron Mass me)

\# 1 MeV approx 1.95695 me

PARTICLE\_DB \= {  
    \# \------------------------------------------------------------------  
    \# TIER 1: FUNDAMENTAL FERMIONS (LEPTONS)  
    \# The geometric ground states.  
    \# \------------------------------------------------------------------  
    'Electron':    {'mass': 1.0,           'tau': 1e30,      'type': 'Lepton'},    \# Stable  
    'Muon':        {'mass': 206.7683,      'tau': 2.197e-6,  'type': 'Lepton'},    \# 105.66 MeV  
    'Tau':         {'mass': 3477.19,       'tau': 2.903e-13, 'type': 'Lepton'},    \# 1776.86 MeV

    \# \------------------------------------------------------------------  
    \# TIER 2: FUNDAMENTAL FERMIONS (QUARKS)  
    \# Current masses (MS-bar) vs Constituent.  
    \# Note: Top is Pole mass. Light quarks are current mass.  
    \# \------------------------------------------------------------------  
    \# NOTE: Quark masses are scale-dependent. Geometric values represent   
    \# the Constituent/Anchor mass. Comparisons use PDG broad ranges.  
    'Up':          {'mass': 4.22,          'tau': 1e30,      'type': 'Quark'},     \# \~2.16 MeV (Anchor Geometry)  
    'Down':        {'mass': 9.13,          'tau': 1e30,      'type': 'Quark'},     \# \~4.67 MeV (Anchor Geometry)  
    'Strange':     {'mass': 186.03,        'tau': 1.2e-8,    'type': 'Quark'},     \# \~95 MeV (MS-bar \~93) \- Used effective  
    'Charm':       {'mass': 2495.0,        'tau': 1.1e-12,   'type': 'Quark'},     \# \~1.27 GeV  
    'Bottom':      {'mass': 8171.5,        'tau': 1.5e-12,   'type': 'Quark'},     \# \~4.18 GeV  
    'Top':         {'mass': 337728.0,      'tau': 5e-25,     'type': 'Quark'},     \# 172.5 GeV (Pole)

    \# \------------------------------------------------------------------  
    \# TIER 3: FUNDAMENTAL BOSONS  
    \# The volumetric connectors.  
    \# \------------------------------------------------------------------  
    'W Boson':     {'mass': 157270.6,      'tau': 3e-25,     'type': 'Boson'},     \# 80.379 GeV  
    'Z Boson':     {'mass': 178449.0,      'tau': 3e-25,     'type': 'Boson'},     \# 91.1876 GeV  
    'Higgs':       {'mass': 245098.7,      'tau': 1.5e-22,   'type': 'Boson'},     \# 125.25 GeV

    \# \------------------------------------------------------------------  
    \# TIER 4: LIGHT MESONS (PSEUDOSCALAR & VECTOR)  
    \# The primary "Noise" in the 200-2000 me range.  
    \# \------------------------------------------------------------------  
    'Pion\_ch':     {'mass': 273.13,        'tau': 2.603e-8,  'type': 'Meson'},     \# 139.57 MeV (pi+/-)  
    'Pion\_0':      {'mass': 264.15,        'tau': 8.4e-17,   'type': 'Meson'},     \# 134.98 MeV  
    'Eta':         {'mass': 1072.1,        'tau': 5.0e-19,   'type': 'Meson'},     \# 547.86 MeV  
    'Rho(770)':    {'mass': 1517.1,        'tau': 4.5e-24,   'type': 'Meson'},     \# 775.26 MeV (Vector)  
    'Omega(782)':  {'mass': 1531.6,        'tau': 7.7e-23,   'type': 'Meson'},     \# 782.65 MeV  
    'Eta\_prime':   {'mass': 1874.3,        'tau': 3.2e-21,   'type': 'Meson'},     \# 957.78 MeV  
    'Phi(1020)':   {'mass': 1995.0,        'tau': 1.5e-22,   'type': 'Meson'},     \# 1019.46 MeV  
    'f0(980)':     {'mass': 1937.0,        'tau': 1e-23,     'type': 'Meson'},     \# \~990 MeV

    \# \------------------------------------------------------------------  
    \# TIER 5: STRANGE MESONS  
    \# \------------------------------------------------------------------  
    'Kaon\_ch':     {'mass': 966.09,        'tau': 1.238e-8,  'type': 'Meson'},     \# 493.67 MeV  
    'Kaon\_0\_S':    {'mass': 973.80,        'tau': 8.9e-11,   'type': 'Meson'},     \# 497.61 MeV  
    'Kaon\_0\_L':    {'mass': 973.80,        'tau': 5.1e-8,    'type': 'Meson'},     \# 497.61 MeV  
    'K\*(892)':     {'mass': 1744.9,        'tau': 1.3e-23,   'type': 'Meson'},     \# 891.66 MeV

    \# \------------------------------------------------------------------  
    \# TIER 6: CHARMED & BOTTOM MESONS  
    \# \------------------------------------------------------------------  
    'D\_0':         {'mass': 3649.4,        'tau': 4.1e-13,   'type': 'Meson'},     \# 1864.84 MeV  
    'D\_ch':        {'mass': 3658.8,        'tau': 1.0e-12,   'type': 'Meson'},     \# 1869.66 MeV  
    'D\_s':         {'mass': 3852.0,        'tau': 5.0e-13,   'type': 'Meson'},     \# 1968.35 MeV  
    'J/Psi':       {'mass': 6060.5,        'tau': 7.1e-21,   'type': 'Meson'},     \# 3096.9 MeV  
    'B\_0':         {'mass': 10332.0,       'tau': 1.5e-12,   'type': 'Meson'},     \# 5279.66 MeV  
    'Upsilon(1S)': {'mass': 18513.3,       'tau': 1.2e-20,   'type': 'Meson'},     \# 9460.3 MeV

    \# \------------------------------------------------------------------  
    \# TIER 7: BARYONS (The Stable & Quasi-Stable Octet)  
    \# \------------------------------------------------------------------  
    'Proton':      {'mass': 1836.15,       'tau': 1e34,      'type': 'Baryon'},    \# 938.27 MeV  
    'Neutron':     {'mass': 1838.68,       'tau': 879.4,     'type': 'Baryon'},    \# 939.56 MeV  
    'Lambda':      {'mass': 2183.3,        'tau': 2.6e-10,   'type': 'Baryon'},    \# 1115.68 MeV  
    'Sigma\_ch':    {'mass': 2327.5,        'tau': 8.0e-11,   'type': 'Baryon'},    \# 1189.37 MeV (+)  
    'Sigma\_0':     {'mass': 2333.9,        'tau': 7.4e-20,   'type': 'Baryon'},    \# 1192.64 MeV  
    'Xi\_0':        {'mass': 2573.1,        'tau': 2.9e-10,   'type': 'Baryon'},    \# 1314.86 MeV  
    'Xi\_ch':       {'mass': 2586.5,        'tau': 1.6e-10,   'type': 'Baryon'},    \# 1321.71 MeV  
    'Omega\_m':     {'mass': 3272.9,        'tau': 8.2e-11,   'type': 'Baryon'},    \# 1672.45 MeV

    \# \------------------------------------------------------------------  
    \# TIER 8: THE GHOST FOREST (Excited Baryonic Resonances)  
    \# These are the short-lived states that create the "Red" noise lines.  
    \# \------------------------------------------------------------------  
    \# The Delta Baryons (Excited Nucleons)  
    'Delta(1232)': {'mass': 2411.0,        'tau': 5.6e-24,   'type': 'Resonance'}, \# 1232 MeV (Width \~117 MeV)  
    'Delta(1600)': {'mass': 3131.0,        'tau': 2.6e-24,   'type': 'Resonance'}, \# \~1600 MeV  
    'Delta(1620)': {'mass': 3190.0,        'tau': 4.5e-24,   'type': 'Resonance'}, \# \~1630 MeV  
    'Delta(1700)': {'mass': 3327.0,        'tau': 2.2e-24,   'type': 'Resonance'}, \# \~1700 MeV

    \# The N Resonances (Roper, etc.)  
    'N(1440)':     {'mass': 2818.0,        'tau': 2.0e-24,   'type': 'Resonance'}, \# 1440 MeV (Roper)  
    'N(1520)':     {'mass': 2974.5,        'tau': 6.0e-24,   'type': 'Resonance'}, \# 1520 MeV  
    'N(1535)':     {'mass': 3004.0,        'tau': 4.3e-24,   'type': 'Resonance'}, \# 1535 MeV  
    'N(1650)':     {'mass': 3229.0,        'tau': 4.3e-24,   'type': 'Resonance'}, \# 1650 MeV

    \# The Excited Hyperons  
    'Sigma(1385)': {'mass': 2710.4,        'tau': 1.8e-23,   'type': 'Resonance'}, \# 1385 MeV  
    'Lambda(1405)':{'mass': 2750.5,        'tau': 1.3e-23,   'type': 'Resonance'}, \# 1405 MeV  
    'Xi(1530)':    {'mass': 2994.0,        'tau': 7.0e-23,   'type': 'Resonance'}, \# 1530 MeV  
    'Lambda\_c':    {'mass': 4474.5,        'tau': 2.0e-13,   'type': 'Resonance'}, \# 2286.46 MeV (Charmed Lambda)  
}  